\subsection{\bt{Enforcement under concurrent writes}{并发写入下的强制保证}}
\label{sec:enforce-eval}
\bi{\emph{Setup (Q2).} We compare checking before execution (\emph{pre-check}, the default) with snapshot-bound execution for store revenue at 200,000 rows. The violating write reloads four months of sales, duplicating their keys. Each probe query returns the canonical answer for one month together with the number of duplicated keys in that month, computed in the same statement and therefore on the same snapshot; an answer is a \emph{violation} if that count is positive, i.e., it was computed on data that violates the declared revision's grain condition.}
{\emph{设置（Q2）。}我们在 20 万行数据上比较执行前检查（\emph{预检查}，默认方式）与绑定快照的执行，指标为门店营业额。违反条件的写入把四个月的销售重新装载一次，使其键重复。每条探针查询返回某个月的规范答案，同时在同一条语句、因而在同一快照上算出该月重复键的个数；该数大于零即为\emph{违规}，即答案是在违反所声明修订粒度条件的数据上算出的。}

\bi{\emph{Deterministic interleavings.} When the write commits after the check and before execution, pre-check answers on violating data in \SnPausePre\ of \SnTrials\ trials and snapshot-bound execution in none, answering from the pre-write snapshot. When the write commits just before an unannounced call, pre-check still answers on violating data in \SnBeforePre\ trials because its version source has not yet observed the write; snapshot-bound execution rejects all \SnBeforeSnapRej. When the write is announced and its statistics are visible, both reject.}
{\emph{确定性交错。}写入在检查之后、执行之前提交时，预检查在 \SnTrials\ 次中有 \SnPausePre\ 次在违规数据上作答，绑定快照的执行一次也没有，它在写入前的快照上作答。写入刚提交、未经通知就调用时，预检查仍有 \SnBeforePre\ 次违规，因为其版本来源尚未观察到这次写入；绑定快照的执行 \SnBeforeSnapRej\ 次全部拒绝。写入经过通知、统计已经可见时，两者都拒绝。}

\begin{table}[t]
\caption{\bt{Random concurrency: four readers issue declaring queries for 3\,s while a writer commits the violating write at a random moment, \SnStressTrials\ trials per row. Readers ask about the same or different months. ``Announced'': the writer logs the batch and notifies the middleware right after commit.}{随机并发：4 个读者在 3 秒内持续发出声明了修订的查询，写者在随机时刻提交违反条件的写入，每行 \SnStressTrials\ 次。读者问相同或不同的月份。“提交后通知”：写者提交后立即登记批次并通知中间层。}}
\label{tab:stress}
\footnotesize
\setlength{\tabcolsep}{3pt}
% generated by tools/review-results.py
\begin{tabularx}{\linewidth}{@{}llLrr@{}}
\toprule
& & \multicolumn{2}{c}{\bt{Pre-check}{预检查}} & \bh{Snapshot-bound}{绑定快照}\\
\cmidrule(lr){3-4}
\bt{Writer}{写者} & \bt{Months}{月份} & \bt{violations / answered}{违规／作答} & \bh{trials}{试验} & \bh{violations / answered}{违规／作答}\\
\midrule
\bt{unannounced}{不通知} & \bt{same}{相同} & 141 / 3{,}077 (4.58\%) & 19 & 0 / 2{,}767\\
\bt{unannounced}{不通知} & \bt{different}{不同} & 120 / 3{,}128 (3.84\%) & 19 & 0 / 2{,}849\\
\bt{announced}{提交后通知} & \bt{same}{相同} & 8 / 2{,}848 (0.28\%) & 2 & 0 / 2{,}818\\
\bt{announced}{提交后通知} & \bt{different}{不同} & 2 / 2{,}934 (0.07\%) & 1 & 0 / 2{,}863\\
\bottomrule
\end{tabularx}

\end{table}


\bi{\emph{Random concurrency.} Table~\ref{tab:stress} randomizes the interleaving. Without announcement, pre-check answers on violating data in \SnPreUnannMin--\SnPreUnannMax\% of answered queries, in at least \SnPreUnannTrials\ of \SnStressTrials\ trials, starting up to \SnUnannLagMax\,ms after the commit, the lag of its version source. Announcing the write immediately after commit narrows but does not close the window: \SnPreAnnViol\ violations remain, all from queries checked before the commit and executed after it. Snapshot-bound execution has no violation in \SnSnapAnswered\ answered queries (95\% upper bound \SnSnapUpper\% per row).}
{\emph{随机并发。}表~\ref{tab:stress} 把交错随机化。不通知时，预检查有 \SnPreUnannMin--\SnPreUnannMax\% 的作答落在违规数据上，\SnStressTrials\ 次试验中至少 \SnPreUnannTrials\ 次出现，违规始于提交后最多 \SnUnannLagMax\ 毫秒内，即其版本来源的滞后。写入提交后立即通知能缩小但不能关闭这个窗口：仍有 \SnPreAnnViol\ 次违规，全部来自提交前已通过核对、提交后才执行的查询。绑定快照的执行在 \SnSnapAnswered\ 次作答中没有违规（每行 95\% 上界 \SnSnapUpper\%）。}

\bi{\emph{Cost.} Under these concurrent readers, both modes have the same median latency (\SnStressPFiftyLo--\SnStressPFiftyHi\,ms). A single reader without writes pays two more round trips (\SnSteadySnap\ vs.\ \SnSteadyPre\,ms); the first query after a write takes \SnFirstSnap\ versus \SnFirstPre\,ms, because the snapshot reuses the verdicts that maintenance recorded under unchanged versions instead of checking twice. Writers pay for the trigger: single-row statements run at \SnWriteSeq$\times$ throughput sequentially and \SnWriteConcTT$\times$ with 32 concurrent writers (\SnWriteHotTT$\times$ with an unsharded counter), while 10,000-row statements take \SnWriteBulk$\times$ as long.}
{\emph{代价。}在上述并发读者下，两种模式的中位延迟相同（\SnStressPFiftyLo--\SnStressPFiftyHi\ 毫秒）。没有写入的单个读者多付两次往返（\SnSteadySnap\ 对 \SnSteadyPre\ 毫秒）；写入后的第一条查询为 \SnFirstSnap\ 对 \SnFirstPre\ 毫秒，因为快照复用了维护在版本未变时记录的结论，而不是再检查一遍。写者为触发器付出代价：单行语句的吞吐在顺序执行时为 \SnWriteSeq\ 倍，32 个并发写者时为 \SnWriteConcTT\ 倍（计数不分片时为 \SnWriteHotTT\ 倍），每条 1 万行的语句耗时为 \SnWriteBulk\ 倍。}

